\section{Questions, Output Contract and Cost}
\label{sec:framework}
We distinguish three questions: can an execution description predict a trajectory, does that trajectory preserve a specified reference, and does a qualified implementation reduce complete cost? Acceptance rate, isolated draft throughput and target-forward batching answer narrower questions.

\subsection{Prediction and reference compatibility}
Let $k$ be maximum draft width, UB the target micro-batch capacity ($n_{\mathrm{ubatch}}$), and EOG a vocabulary end-of-generation ID. A greedy speculative round verifies proposals, commits the agreeing prefix and a correction or bonus when budget permits, then removes inconsistent suffix state. Our strict output contract compares complete committed generated IDs with the matched target-only greedy reference under the frozen model/runtime/configuration. A failed pair blocks a workload-global lossless optimization claim; a small passing probe establishes eligibility for a larger gate. Capped ID equality does not measure functional HumanEval correctness or pass@$k$.

Historical text-path agreement (H) compares arrays reconstructed from CLI text. Direct-ID qualification (Q) compares captured IDs with the study's reference. Prediction equality compares a sealed complete array with reserved native live IDs. Prediction can succeed on a reference-divergent path, so it does not imply Q. We write prediction equality in words; the letter P below denotes only the draft worker. Same-state intervention has separate score, immediate-decision and continuation endpoints.

\subsection{Execution descriptions as testable candidates}
Table~\ref{tab:inputcontracts} states what each program receives and computes. UB denotes capacity, not a trace of actual forward shapes, occupied positions or prior KV operations. In particular, target-only UB2 and native speculative UB2 need not issue the same sequence of forwards. The candidate test checks complete output directly. Equal parameter names alone do not establish equal numerical paths.
\begin{table*}[t]
\caption{Inputs and local execution contracts. Each program starts from its own prompt prefill. P = draft worker; R = target worker. Shared numerical routines limit implementation independence. UB labels do not specify identical actual forward traces.}
\label{tab:inputcontracts}
\centering
\begin{tabular}{@{}>{\raggedright\arraybackslash}p{1.03in}>{\raggedright\arraybackslash}p{1.43in}>{\raggedright\arraybackslash}p{2.00in}>{\raggedright\arraybackslash}p{1.54in}@{}}
\toprule
Program & Generation inputs & Locally computed history / decisions & Separation or missing equivalence \\
\midrule
B570 H1 & Prompt; target weights; UB1; no draft or proposals & Own target-only generated prefix, KV and greedy decisions & Draft-free predictor; no unique-cause isolation \\
RTX target-only UB2 & Prompt; target weights; UB2; no draft or proposals & Own target-only prefix and KV; one generated-ID stream & Capacity-level candidate; equality of actual native forward traces is not established \\
RTX H2 & Prompt; P's proposals from its own earlier R feedback & R checks own prefix/pending ID; computes sampling, accept/correct/bonus and KV rollback & Separate P/R orchestration; shared llama.cpp/GGML/CUDA \\
RTX fresh R\_ONLY & Prompt; sealed H2 proposal requests, read sequentially & Fresh R computes its own target decisions and state transitions & Target-only execution of tape; tape generation uses draft; same numerical routines \\
Native RTX live & Prompt; its own draft proposals; no tape import & Native target/draft history, acceptance and state transitions & Independent live outcome; not an independent numerical implementation \\
\bottomrule
\end{tabular}
\end{table*}


\subsection{Complete realized cost}
Let $T_b$ be matched target-only cost per output token, $T_d^{\mathrm{real}}$ complete realized draft-subsystem time including interface replay, $T_{\mathrm{target,round}}^{\mathrm{real}}$ complete target work, and $T_x$ residual synchronization, transfer and bookkeeping outside those regions. Let $C$ count all committed output per round, including correction/bonus. Then
\begin{align}
S &= \frac{T_b C}{T_d^{\mathrm{real}}+T_{\mathrm{target,round}}^{\mathrm{real}}+T_x}\label{eq:physical}\\
 &= \frac{C}{\dreal+\rtarget+X}.\label{eq:normalized}
\end{align}
The normalized denominator gives break-even $C>\dreal+\rtarget+X$. These are accounting identities, not new profitability theory \citep{specllm}. Total-time ratios need not equal prompt geometric means. Incompatible experiments cannot isolate an unmeasured component by subtraction. The historical target-forward proxy $\microv$ omits drafting and integration; its original cells and separate fresh qualification gates remain in Appendix F.

Primary results retain their frozen questions and gates; post-study results are identified separately. Timing on blocked paths remains diagnostic, and a later Q gate does not certify historical H timing.
